\documentclass[10pt,a4paper]{amsart}
\usepackage[margin=19mm]{geometry}
\usepackage{amsmath,amssymb,mathtools}
\usepackage[scr=rsfs,cal=euler]{mathalfa}
\usepackage{xfrac}
\usepackage[unicode,hidelinks,bookmarks=false]{hyperref}
\newcommand{\ie}{i.e.\ }
\newcommand{\cf}{cf.\ }
\newcommand{\resp}{respectively\ }
\newcommand{\Subsection}{\subsection}
\providecommand{\nobreakdash}{\nobreak\hbox{-}}
\setlength{\emergencystretch}{2em}
\multlinegap0pt
\usepackage{amssymb}
\usepackage{mathtools}
\usepackage{algorithm}\usepackage{algpseudocode}
\usepackage{xcolor}
\newtheorem{prop}{Proposition}
\newcommand{\Mat}[1]{#1}
\newcommand{\tH}{\widetilde{H}}
\newcommand{\tS}{\widetilde{S}}
\title{A scalable preconditioning framework for stabilized contact mechanics with hydraulically active fractures}
\author{Andrea Franceschini, Laura Gazzola, Massimiliano Ferronato}
\date{Source: arXiv:2112.12397v1 (CC BY 4.0)}
\begin{document}
\maketitle

\noindent\emph{Source and licence.} A scalable preconditioning framework for stabilized contact mechanics with hydraulically active fractures. Version arXiv:2112.12397v1 (CC BY 4.0), distributed by arXiv under the Creative Commons Attribution 4.0 International licence, http://creativecommons.org/licenses/by/4.0/, as recorded in the arXiv licence metadata for this exact version and shown on the abstract page https://arxiv.org/abs/2112.12397v1. Full licence notice: \texttt{licenses/arxiv-2112.12397-CC-BY-4.0.txt}.\par\smallskip

\noindent\textit{Editorial scope.} Counted unit: the article's prop thm:eig\_tup (source0.tex lines 1179--1199). The article's complete original proof of it is delivered with it (source0.tex lines 1201--1225, one \texttt{proof} environment of 25 lines). No mathematics is written, repaired or re-derived: the statement and the proof are the article's own lines, in the article's own notation, and no retained line is substituted or re-flowed.  Retained uncounted as the source-local context the proof needs: lines 851--854; lines 1083--1137; lines 1169--1177.  Nothing was omitted from any retained span, so the package carries no omission record.  No retained line contains a citation, so the wrapper carries no bibliography and no bibliography entry.  No retained line contains a figure environment, an image-inclusion command or drawing code, so no image is delivered and the package declares no asset dependency.  Editorial formatting only, in addition: an AMS \texttt{amsart} preamble that restates the article's own declarations and macros used by the retained lines, namely the environments \texttt{prop} on separate counters, together with the standard packages those lines need.  The retained environments print in this document as Proposition 1 in order of appearance, whereas the article prints the same items with the numbering of its own full document.  The complete native source of the article as distributed in the arXiv e-print for this exact version is bundled unchanged as \texttt{sources/120046/source0.tex}.
\begin{align}
    E_H &= I_t - \tH^{-1} H, \label{eq:EH} \\
    E_S &= I_u - \tS^{-1} S, \label{eq:ES}
\end{align}

Similarly to the t-p-u approach, the preconditioning operator $\mathcal{M}_2^{-1}$ can be written as an inexact block LDU factorization of $\mathcal{J}$. With
the permutation matrix $\mathcal{Q}_2$:
\begin{equation}
  \mathcal{Q}_2 = \begin{bmatrix}
    \Mat{0} & \Mat{I}_t & \Mat{0} \\
    \Mat{I}_u & \Mat{0} & \Mat{0} \\
    \Mat{0} & \Mat{0} & \Mat{I}_p \\
  \end{bmatrix},
  \label{eq:Q2}
\end{equation}
the block LDU factorization reads:
\begin{equation}
  \mathcal{Q}_2 \mathcal{J} \mathcal{Q}_2^T \simeq \mathcal{L}_2 \mathcal{D}_2 \mathcal{U}_2,
\end{equation}
with:
\begin{align}
  \mathcal{L}_2 &= \begin{bmatrix}
    \Mat{I}_t & \Mat{0} & \Mat{0} \\
    -\Mat{C}_1\Mat{\tH}^{-1} & \Mat{I}_u & \Mat{0} \\
    \Mat{0} & \Mat{Q}_2\Mat{\tS}_1^{-1} & \Mat{I}_p \\
  \end{bmatrix}, &&&
  \mathcal{D}_2 &= \begin{bmatrix}
    -\Mat{\tH} & \Mat{0} & \Mat{0} \\
    \Mat{0} & \Mat{S}_1 & \Mat{0} \\
    \Mat{0} & \Mat{0} & \Mat{\tS}_2 \\
  \end{bmatrix} & \text{and} &&
  \mathcal{U}_2 &= \begin{bmatrix}
    \Mat{I}_t & -\Mat{\tH}^{-1}\Mat{C}_2 & \Mat{0} \\
    \Mat{0} & \Mat{I}_u & \Mat{\tS}_1^{-1}\Mat{Q}_1 \\
    \Mat{0} & \Mat{0} & \Mat{I}_p \\
  \end{bmatrix}.
  \label{eq:LDU2}
\end{align}
%where $\Mat{S}_1$ and $\Mat{S}_2$ are the two Schur complements, defined as:
%\begin{align}
%  \Mat{S}_1 &= \Mat{A} - \Mat{C}_1 \Mat{\tH}^{-1} \Mat{C}_2, \label{eq:S1} \\
%  \Mat{S}_2 &= \Mat{T} - \Mat{Q}_2 \Mat{S}_1^{-1} \Mat{Q}_1. \label{eq:S2} 
%\end{align}
%As before, $\Mat{S}_1$ is a symmetric negative semidefinite matrix when all the elements are in stick mode. 
The final algebraic expression of $\mathcal{M}_2^{-1}$ is therefore:
\begin{equation}
    \mathcal{M}_2^{-1} = \left[ \begin{array}{ccc}
    I_t & \tH^{-1}C_2 & -\tH^{-1} C_2 \tS_1^{-1} Q_1 \\
    0 & I_u & -\tS_1^{-1} Q_1 \\
    0 & 0 & I_p \end{array} \right]
    \left[ \begin{array}{ccc}
    -\tH^{-1} & 0 & 0 \\
    0 & \tS_1^{-1} & 0 \\
    0 & 0 & \tS_2^{-1} \end{array} \right]
    \left[ \begin{array}{ccc}
    I_t & 0 & 0 \\
    C_1 \tH^{-1} & I_u & 0 \\
    -Q_2 \tS_1^{-1} C_1 \tH^{-1} & -Q_2 \tS_1^{-1} & I_p \end{array} \right].
    \label{eq:P2}
\end{equation}

Let us introduce the matrices:
\begin{subequations}
\begin{align}
    E_{S_1} &= I_u - \tS_1^{-1} S_1, \label{eq:ES1} \\
    E_{S_2} &= I_p - \tS_2^{-1} \hat{S}_2, \label{eq:ES2}
\end{align}
\label{eq:E2def}\null
\end{subequations}
where $\hat{S}_2=T-Q_2\tS_1^{-1}Q_1$. Along with $E_H$ already introduced in equation \eqref{eq:EH}, $E_{S_1}$ and $E_{S_2}$ are matrix measures of the quality of the approximations $\tS_1$ and $\tS_2$ introduced in $\mathcal{M}_2^{-1}$. The following result holds.

\begin{prop}
\label{thm:eig_tup}
The eigenvalues $\lambda_2\in\mathbb{C}$ of the preconditioned matrix with the t-u-p approach are such that:
\begin{equation}
    |\lambda_2-1| \leq (1+\gamma) \max\{\varepsilon_H,\varepsilon_{S_1}, \varepsilon_{S_2}\},
    \label{eq:eigbound2}
\end{equation}
with $\varepsilon_H=\|E_H\|$, $\varepsilon_{S_1}=\|E_{S_1}\|$, $\varepsilon_{S_2}=\|E_{S_2}\|$, $\gamma=\|\mathcal{G}\|$ and
\begin{equation}
    \mathcal{G} = \left[ \begin{array}{ccc}
    C_{2H}\hat{I}_QC_{1S} & -C_{2H}\hat{I}_Q & C_{2H}Q_{1S} \\
    \hat{I}_QC_{1S} & -Q_{1S}Q_{2S} & Q_{1S} \\
    -Q_{2S}C_{1S} & Q_{2S} & 0 \end{array} \right], 
    \label{eq:Gdef}
\end{equation}
\begin{equation}
    C_{1S} = \tS_1^{-1}C_1, \quad C_{2H} = \tH^{-1}C_2, \quad Q_{1S} = \tS_1^{-1}Q_1, \quad Q_{2S} = \tS_2^{-1}Q_2, \quad \hat{I}_Q = I_u + Q_{1S}Q_{2S},
    \label{eq:matdef}
\end{equation}
for any compatible matrix norm.
\end{prop}

\begin{proof}
Recalling equations \eqref{eq:Q2} and \eqref{eq:P2} and introducing the error matrices \eqref{eq:EH} and \eqref{eq:E2def}, the preconditioned matrix with the t-u-p approach reads:
\small
\begin{equation}
    \mathcal{M}_2^{-1} \mathcal{Q}_2 \mathcal{J} \mathcal{Q}_2^T = \left[ \begin{array}{ccc}
    I_t + \left[ -I_t+\tH^{-1}C_2 \left( I_u + \tS_1^{-1}Q_1\tS_2^{-1}Q_2 \right) \tS_1^{-1}C_1 \right] E_H & -\tH^{-1}C_2 \left( I_u + \tS_1^{-1}Q_1\tS_2^{-1}Q_2 \right) E_{S_1} & \tH^{-1}C_2\tS_1^{-1}Q_1 E_{S_2} \\
    \left( I_u + \tS_1^{-1}Q_1\tS_2^{-1}Q_2 \right) \tS_1^{-1}C_1 E_H & I_u - \left( I_u + \tS_1^{-1}Q_1\tS_2^{-1}Q_2 \right) E_{S_1} & \tS_1^{-1}Q_1 E_{S_2} \\
    -\tS_2^{-1}Q_2\tS_1^{-1}C_1 E_H & \tS_2^{-1}Q_2 E_{S_1} & I_p-E_{S_2} \end{array} \right].
    \label{eq:precmat_tup}
\end{equation}
\normalsize
By making use of the definitions \eqref{eq:Gdef} and \eqref{eq:matdef}, equation \eqref{eq:precmat_tup} can be re-written as:
\begin{equation}
    \mathcal{M}_2^{-1} \mathcal{Q}_2 \mathcal{J} \mathcal{Q}_2^T = \mathcal{I} + \left( \mathcal{G} - \mathcal{I} \right) \mathcal{E}_2,
    \label{eq:precmat_tup2}
\end{equation}
with:
\begin{equation}
    \mathcal{I} = \left[ \begin{array}{ccc}
    I_t & 0 & 0 \\ 0 & I_u & 0 \\ 0 & 0 & I_p \end{array} \right], \qquad \mathcal{E}_2 = \left[ \begin{array}{ccc}
    E_H & 0 & 0 \\ 0 & E_{S_1} & 0 \\ 0 & 0 & E_{S_2} \end{array} \right].
    \label{eq:I_E2}
\end{equation}
The eigenvalues $\lambda_2$ of the matrix in equation \eqref{eq:precmat_tup2} satisfy the bound \eqref{eq:eigbound2}, thus closing the proof.
\end{proof}

\end{document}
